% =====================================================================
% Chapter 1: Introduction
% =====================================================================
\chapter{Introduction}
\label{ch:introduction}

\section{Background}
A startup incubation centre supports early-stage companies for a fixed period. A typical
centre carries out the following activities:
\begin{enumerate}
  \item keeping records of registered startups and their founders;
  \item receiving applications for its incubation programme;
  \item having expert evaluators score each application;
  \item admitting the best applicants into incubation;
  \item assigning mentors to admitted startups; and
  \item tracking each admitted startup's progress through milestones.
\end{enumerate}
Each activity produces data that depends on the data of earlier activities: an evaluation
belongs to an application, an incubation is created from an accepted application, and
mentors and milestones belong to an incubation. This chain of dependent data makes the domain
a good fit for a relational database.

\section{Problem Statement}
\label{sec:problem}
Small incubation centres frequently manage this process with spreadsheets, e-mail threads and
paper forms. This causes four concrete problems:
\begin{itemize}
  \item \textbf{Redundancy and update anomalies.} A startup's details are re-typed in several
        sheets (applications, evaluations, mentoring). When one copy is changed, the others
        become inconsistent.
  \item \textbf{No enforcement of business rules.} Nothing prevents a mentor from being
        ``assigned'' to a startup that was never admitted, or a score of 150 out of 100 being
        recorded.
  \item \textbf{Broken references.} An evaluation can refer to an application that has been
        deleted or never existed.
  \item \textbf{Slow reporting.} Questions such as ``What is the average score of each
        applicant?'' or ``Which mentor guides the most startups?'' must be answered by manual
        counting.
\end{itemize}

\section{Need for the System}
A relational database with an application on top of it addresses each of these problems:
\begin{itemize}
  \item \textbf{Single source of truth:} each fact is stored once, in a normalized schema.
  \item \textbf{Integrity enforced by the database itself:} primary keys, foreign keys,
        \code{NOT NULL}, \code{UNIQUE} and \code{CHECK} constraints reject invalid data even if
        the user interface has a defect.
  \item \textbf{Consistent multi-step operations:} transactions guarantee that, for example,
        accepting an application and creating its incubation record either both happen or
        neither happens.
  \item \textbf{Instant reports:} joins, aggregate functions and subqueries answer management
        questions directly from the stored data.
\end{itemize}

\section{Objectives}
\label{sec:objectives}
The objectives of the project are:
\begin{enumerate}
  \item To design a normalized relational database (at least 3NF) for the incubation process.
  \item To enforce data integrity using appropriate database constraints.
  \item To provide a working user interface through which users perform \code{INSERT},
        \code{UPDATE}, \code{DELETE} and \code{SELECT} operations in every module.
  \item To implement the application life cycle: submission $\rightarrow$ evaluation
        $\rightarrow$ decision $\rightarrow$ incubation.
  \item To support mentor assignment and milestone-based progress tracking for incubated
        startups.
  \item To generate meaningful reports using joins, aggregates, \code{GROUP BY}/\code{HAVING},
        subqueries and views.
  \item To demonstrate transactions for multi-step operations.
\end{enumerate}

\section{Scope of the Project}
\label{sec:scope}
\textbf{In scope.} Five functional modules (Startup, Application, Evaluation, Mentor and
Progress Management), the reports listed in Section~\ref{sec:reports}, the constraints implied
by the business rules, a role-based user interface, a locally hosted PostgreSQL database and
realistic sample data (10 startups, 18 founders, 10 applications, 20 evaluations, 5 evaluators,
6 mentors, 5 incubations and 15 milestones).

\textbf{Out of scope.} The following were deliberately excluded to keep the project focused on
database design and implementation: user login and authentication; payments, funding rounds
and investor portals; notifications, e-mail and chat; AI-based recommendations; multiple
incubation centres; resource or office-space allocation; file uploads; and cloud deployment.

\section{Users of the System}
\label{sec:users}
The application offers a separate set of screens (tabs) for each of four roles,
summarised in Table~\ref{tab:users}.

\begin{table}[H]
  \centering
  \caption{Users of the system and what each role can do}
  \label{tab:users}
  \begin{tabularx}{\textwidth}{@{}lX@{}}
    \toprule
    \textbf{Role} & \textbf{Capabilities} \\
    \midrule
    Administrator      & Register and maintain startups and founders; maintain evaluators and
                         mentors; change application status; admit startups to incubation;
                         assign mentors; create milestones; view all reports. \\
    Startup / Founder  & View and update its own startup; submit applications; track
                         application status and read evaluator feedback; view its milestones
                         and progress. \\
    Evaluator          & See the applications under review that he or she has not yet
                         evaluated; record a score and feedback; view past evaluations. \\
    Mentor             & See the startups he or she mentors; add mentorship remarks; update
                         the milestone status of those startups only. \\
    \bottomrule
  \end{tabularx}
\end{table}

The role is chosen from a ``View as'' selector in the sidebar, after which the user selects
which startup, evaluator or mentor he or she is. Every query on those screens is then filtered
by the selected identifier (for example \code{WHERE evaluator\_id = \%s}). This is a user
interface feature that shows different users different parts of the same database; it is
\emph{not} an authentication mechanism, which is outside the scope of the project.
